\section{Introduction}
\label{sec:intro}

Differential privacy~\cite{dwork2006calibrating} is widely regarded as the
canonical formal guarantee for privacy-preserving machine learning. In its
classical form, a single privacy budget $(\eps, \delta)$ is shared by every
record in the database, and every record is protected to exactly the same
extent. Yet privacy preferences are heterogeneous: a clinical patient and a
casual user contributing data to the same model often place very different
values on their privacy. Recognising this, \emph{individual differential
privacy} (iDP) replaces the single $(\eps,\delta)$ guarantee with a per-record
contract $(\eps_i,\delta)$, and a number of mechanisms have been proposed to
deliver such per-record contracts in deep
learning~\cite{boenisch2023have,jorgensen2015conservative,feldman2021individual}.
The most widely-adopted realisation, sampling-based iDP introduced by
Boenisch et al.~\cite{boenisch2023have}, achieves heterogeneous budgets by
calibrating each record's Poisson inclusion probability $p_i$ against a single
shared noise multiplier $\sigma$ so that, for fixed batch size and iteration
count, every group of records satisfies its own $(\eps_g, \delta)$-DP bound.

We focus on a fundamental and overlooked feature of this realisation.
\emph{Because every record shares the same noise multiplier and a common
expected batch size, each record's sampling rate is mechanically coupled to
every other record's chosen budget.} Concretely, lowering one's own $\eps_i$
forces other contributors' sampling rates upward to keep the expected batch size
fixed---which, as recent empirical work has shown~\cite{kaiser2026your},
inflates their membership-inference (MIA) advantage even though every
formal $(\eps_g,\delta)$-iDP guarantee continues to hold. The promise of
individual control is therefore broken in a structural sense: an individual's
\emph{actual} privacy risk depends not on her own choice but on the joint
choice of the entire population.

\paragraph{Our viewpoint: this is a game.}
The natural framing for any choice problem whose outcome is determined by the
choices of multiple self-interested parties is game theory. Pej\'o et
al.~\cite{pejo2019together} introduced a two-player game on the
privacy-accuracy plane (the \emph{Price of Privacy}), but they were studying a
mechanism in which a player's privacy choice affected only her \emph{own}
accuracy. The novel feature of sampling-based iDP is that every player's
privacy choice affects every other player's \emph{privacy}, a far stronger and
qualitatively different coupling. We call this the \emph{privacy externality},
and the game it gives rise to the \emph{Privacy Externality Game} (PEG).

\paragraph{Contributions.}
This paper makes the following contributions.
\begin{enumerate}
\item \textbf{Formalisation.} We define the $n$-player Privacy Externality
Game (PEG) in which each player picks a per-group budget $\eps_i$ and the
sampling-based iDP mechanism then determines the noise multiplier $\sigma$ and
the per-group sampling rates $p_g$. We split player utility into two awareness
regimes that capture realistic contracting situations: \emph{naive}
players, who price privacy via their nominal $\eps_i$, and
\emph{aware} players, who price privacy via the realised MIA advantage
$A_i(\boldsymbol{\eps})$.
\item \textbf{Externality lemma.} We prove that the externality is monotone:
under the binding constraints of the mechanism, $\partial A_i / \partial \eps_j
\leq 0$ for every $j \neq i$. Lowering any other player's budget never reduces
my realised vulnerability.
\item \textbf{Equilibrium characterisation.} For the aware regime, we observe
that best-response dynamics typically cycle; we therefore characterise the
empirical mixed equilibrium via fictitious play. For the naive regime, every
player's incentive is independent (modulo the global model-utility term), and a
pure Nash equilibrium exists; it coincides with the regime that maximally
exposes all players, since naive players ignore the externality entirely.
\item \textbf{Coalition Stackelberg.} We model a colluding subset of data
holders as a Stackelberg leader maximising a target's realised MIA advantage.
We show this exactly recovers the budget-manipulation attack of Kaiser et
al.~\cite{kaiser2026your} as a one-step best response.
\item \textbf{Empirical evaluation.} We numerically evaluate the game in a
range of small but representative scenarios that mirror Kaiser et al.'s
empirical setup. We report the equilibrium budgets, the equilibrium MIA
advantage, and the realised externality across coalition sizes.
\end{enumerate}

\paragraph{Implication.}
The previous attack literature on iDP frames excess vulnerability as a
\emph{vulnerability}---something a malicious party exploits. We argue the
opposite. Under sampling-based iDP, the externality is the equilibrium
behaviour of rational players, malicious or not. Mitigation therefore cannot
proceed by detecting attackers; it must proceed by changing the contract so
that the underlying game ceases to admit harmful equilibria.

The paper is organised as follows. \Cref{sec:background} reviews the necessary
notions from differential privacy, the sampling-based iDP mechanism, MIA
advantage, and game theory. \Cref{sec:game} formalises the Privacy
Externality Game. \Cref{sec:analysis} proves the externality lemma and
characterises equilibrium properties.  \Cref{sec:experiments} reports the
numerical evaluation. \Cref{sec:related} reviews related work and
\cref{sec:discussion} concludes.
